%============================================================ 초록
%  압축판. 방법 세부·다중비교·효과크기 순위는 본문 몫이라 두지 않는다. 수치는 영문 초록과 집합이 같아야 한다.
\begin{abstract}
\noindent
저비용 자폭 무인수상정(USV) 군집은 전자전이 유선 조종 표적에 닿지 않고 요격탄은 표적보다
비싸 포화되는, 재래식 대응의 사각지대에 있다. 본 연구는 적보다 느린 소수의 방어정이 모선의
접근 회랑을 그물벽으로 선제 차단하는 두 계층 방어 체계를 제안한다. 전략 계층에서는 대규모
언어모델(LLM) 지휘관이 방어정별 담당 무리를 정하고, 기동 계층에서는 U-Net 이 전장 격자
지도 위에 낸 점수맵에서 각 방어정이 픽셀을 지목해 그물벽 위치를 정한다. 두 계층은 같은
주기로 갱신되지만 응답 지연이 크게 달라 느린 전략 계층만 비동기로 분리한다. 픽셀 지목은
행동공간을 격자에 가두어 안전·임무 제약을 마스킹으로 직접 부과하고 적선 수에 불변이며,
학습은 가치망 없는 그룹 상대 정책 최적화를 다중 에이전트로 확장한다. 조합 최적화 후처리를
두지 않은 조건에서 8개 조건 $\times$ 3개 포메이션 $\times$ 30개 시드, 720 에피소드를
시드 대응표본으로 수집해 두 계층의 기여를 $2\times2$ 로 분리 측정하였다. 기동 계층은 LLM
없이 단독으로 포획률을 $0.844 \to 0.886$ 으로 높였고, 제안 체계 전체는 같은 방어 수준에서
포획당 그물 소모를 44.5\,\%, 총 선회량을 46.6\,\% 줄였으며, 두 계층의 기여는 상승이 아닌
\emph{가산적}이었다. 지휘관은 능력 임계 위에서만 이득이었다 --- 클라우드 모델은 $+0.060$,
로컬 Qwen2.5 7B 는 $-0.192$ 로 코드 휴리스틱보다 유의하게 나빴다. 지휘관 성능을 설명한 것은
배정 커버리지가 아니라 재계획 간 계획 안정성으로, 주기적으로 재계획하는 구조에서는 국소
최적성보다 시간에 걸친 일관성이 중요하다는 설계 지침을 얻었다.
\end{abstract}
